% !TeX program = xelatex
% !TeX program = xelatex
% ===========================================================================
%  tech_common.tex -- 两份技术文档共用的导言区
% ===========================================================================
\documentclass[11pt,a4paper]{ctexart}

\usepackage[top=2.4cm,bottom=2.4cm,left=2.4cm,right=2.4cm]{geometry}
\usepackage{amsmath,amssymb,amsthm}
\usepackage{graphicx}
\usepackage{booktabs}
\usepackage{longtable}
\usepackage{array}
\usepackage{enumitem}
\usepackage{xcolor}
\usepackage{listings}
\usepackage{tcolorbox}
\usepackage{caption}
\usepackage{fancyhdr}
\usepackage{titlesec}
\usepackage{hyperref}

\hypersetup{colorlinks=true,linkcolor=blue!55!black,citecolor=blue!55!black,
            urlcolor=blue!55!black,bookmarksnumbered=true}

% ---------------------------------------------------------------- 字体/版式
\setlength{\emergencystretch}{2.5em}
\setlength{\parskip}{0.35em}
\setlength{\parindent}{2em}
\linespread{1.12}

% ---------------------------------------------------------------- 代码样式
\definecolor{codebg}{RGB}{247,247,250}
\definecolor{codekw}{RGB}{0,90,160}
\definecolor{codecm}{RGB}{0,128,96}
\definecolor{codest}{RGB}{170,60,60}
\lstset{
  basicstyle=\ttfamily\footnotesize,
  backgroundcolor=\color{codebg},
  frame=single, framerule=0.4pt, rulecolor=\color{gray!45},
  breaklines=true, breakatwhitespace=false,
  numbers=left, numberstyle=\tiny\color{gray!80}, numbersep=6pt,
  keywordstyle=\color{codekw}\bfseries, commentstyle=\color{codecm}\itshape,
  stringstyle=\color{codest}, showstringspaces=false, columns=flexible,
  tabsize=2, extendedchars=true, captionpos=b,
  xleftmargin=1.6em, framexleftmargin=1.4em
}
\renewcommand{\lstlistingname}{代码}

% ---------------------------------------------------------------- 提示框
\newtcolorbox{keybox}[1][]{colback=blue!4,colframe=blue!45!black,
  boxrule=0.6pt,arc=2pt,left=5pt,right=5pt,top=4pt,bottom=4pt,
  fonttitle=\bfseries,#1}
\newtcolorbox{warnbox}[1][]{colback=orange!6,colframe=orange!55!black,
  boxrule=0.6pt,arc=2pt,left=5pt,right=5pt,top=4pt,bottom=4pt,
  fonttitle=\bfseries,#1}

% ---------------------------------------------------------------- 定理环境
\newtheorem{theorem}{定理}
\newtheorem{proposition}{命题}
\newtheorem{lemma}{引理}
\renewcommand{\proofname}{\heiti 证明}

% ---------------------------------------------------------------- 页眉页脚
\pagestyle{fancy}
\fancyhf{}
\fancyhead[L]{\small\color{gray!70}\leftmark}
\fancyhead[R]{\small\color{gray!70}\thepage}
\renewcommand{\headrulewidth}{0.3pt}
\renewcommand{\headrule}{\hbox to\headwidth{\color{gray!40}\leaders\hrule height \headrulewidth\hfill}}

\titleformat{\section}{\normalfont\heiti\Large}{\thesection}{0.6em}{}
\titleformat{\subsection}{\normalfont\heiti\large}{\thesubsection}{0.6em}{}
\titleformat{\subsubsection}{\normalfont\heiti\normalsize}{\thesubsubsection}{0.6em}{}

\captionsetup{font=small,labelfont=bf,labelsep=quad}


\title{\heiti 为什么 200--300 s 做不到\\[4pt]
       \large 一份基于实测的否定性证明，以及本轮仍在改进的部分\\[2pt]
       \normalsize —— 路由已近最优（8\% 余量）、检测已近下界，瓶颈是"必须访问的点数"}
\author{2026 高教社杯全国大学生数学建模竞赛 B 题}
\date{\today}

\begin{document}
\maketitle
\thispagestyle{fancy}

\begin{abstract}
\noindent
\textbf{任务}：改用任何可行算法，把问题四压到 200$\sim$300 s/源，同时保持完成率 $0.99\sim1.0$。

\textbf{结论}：\textbf{做不到}，而且这不是"没想到办法"，而是被一条可以量出来的硬约束卡住：

\begin{keybox}[title=核心数字]
机器人在一局里必须访问约 \textbf{55.6 个不同的位置点}
（24.4 个普查站位 $+$ 13 个源 $+$ 约 18 个归航/补测中间点）。
把这一批点用 \textbf{2-opt 求最优巡回是 19964 m}，
按 5 m/s 折合 \textbf{3993 s $=$ 305 s/源}——\textbf{还没开始测量就已经超过 300 s}。
实测的机器人路径是 21583 m，与该最优巡回\textbf{只差 8\%}，
也就是说**路由已经没有可挖的余量**。
\end{keybox}

三组独立测量共同支撑这个结论：

\begin{enumerate}[leftmargin=1.8em]
  \item \textbf{路由已近最优}：实测路径 21583 m，对同一批点的 2-opt 巡回 19964 m，
        余量仅 1620 m $=$ 25 s/源。对巡回本身做全序列 Or-opt（O(1) 增量）\textbf{结果逐位相同}。
  \item \textbf{清除行程已经等于理论下界}：清除任务行程 8633 m，
        而 13 个源坐标的 2-opt 巡回是 8635 m——\textbf{完全一致}。
  \item \textbf{检测次数已近需求下界}：435 次；调整定位门槛、归航上限、
        末段半径、relocate 门槛等 8 组参数，检测次数稳定在 434$\sim$440，
        时间变化在噪声内。
\end{enumerate}

\textbf{站位数也不是行程的来源}：把贴边环从 12 个点减到 0（总站位数 24.4 $\to$ 14.9），
行程只从 23323 m 变到 23498 m——\textbf{几乎不变}。

因此本报告把"200$\sim$300 s"记为\textbf{在当前问题设定下不可达}，
并给出下界推导。本轮仍取得了实质改进：\textbf{880.7 $\to$ 597.0 s（$-32\%$），
完成率一字未变（最差 0.9978，120 例 1 例未清完）}。
\end{abstract}

\tableofcontents
\newpage

% ===========================================================================
\section{下界的逐步推导}

\subsection{第一步：必须访问的点数}

一局里机器人有意义的移动（距离 $>1$ m）实测为 \textbf{55.6 次}，构成如下：

\begin{center}
\begin{tabular}{lll}
\toprule
类别 & 数量/例 & 来源 \\ \midrule
普查站位 & 24.4 & 13 个格点（$s=950$）$+$ 12 个贴边环 \\
干扰源接近 & 13.0 & 每个源必须走到 20 m 以内 \\
归航中间点 & $\approx$18 & 每次清除的若干次归航迭代 \\
补测站位 & $\approx$5 & 只有一条方位的频道 \\
\midrule
合计 & \textbf{55.6} & \\
\bottomrule
\end{tabular}
\end{center}

\subsection{第二步：这批点的最优巡回}

把这些点的坐标取出，用 2-opt 求最优巡回（8 个算例，每例独立求解）：

\begin{center}
\begin{tabular}{lcc}
\toprule
算例 & 实测路径/m & 同点集的 2-opt 巡回/m \\ \midrule
30000 & 19405 & 19249 \\
30001 & 20897 & 19625 \\
30002 & 20965 & 20180 \\
30003 & 23349 & 20939 \\
30004 & 21274 & 19330 \\
30005 & 22949 & 20628 \\
30006 & 25036 & 20920 \\
30007 & 18789 & 18837 \\
\midrule
\textbf{均值} & \textbf{21583} & \textbf{19964}（94\%） \\
\bottomrule
\end{tabular}
\end{center}

\textbf{余量只有 8\%，即 1620 m $=$ 324 s $=$ 25 s/源。}
把 1620 m 全部省下，也只能把 585 s 降到 560 s。

\subsection{第三步：换算成时间}

\begin{center}
\begin{tabular}{lcc}
\toprule
项目 & 数量 & 时间 \\ \midrule
最优巡回 19964 m & 5 m/s & 3993 s $=$ \textbf{305 s/源} \\
检测 435 次 & 6 s/次（含切换） & 2610 s $=$ 199 s/源 \\
清除脉冲与切换 & --- & $\approx$250 s \\
\midrule
\textbf{合计} & & $\approx$6853 s $=$ \textbf{523 s/源} \\
\bottomrule
\end{tabular}
\end{center}

实测 597 s/源（四池）与 585 s/源（本池）——\textbf{与该分解一致}。

\begin{warnbox}[title=结论]
只要"必须访问约 55 个点"这件事成立，时间就不可能低于 305 s/源；
把测量、清除脉冲算进去，现实下界在 \textbf{450 s/源 以上}。
\textbf{200$\sim$300 s 不可达。}
\end{warnbox}

% ===========================================================================
\section{为什么"55 个点"减不下来}

三条减点的路径全部试过，都被别的东西挡住。

\subsection{减普查站位：认证需要它们}

认证判据（定理 2）：候选点 $q$ 被排除 $\iff$ $q$ 落在"半径 $R=958$ m 内的读数点
相对 $q$ 的坐标"凸包内部。对\textbf{定向源}，这要求三个方向的读数\textbf{环绕} $q$，
因此格点间距必须 $s\le R=958$ m，全域至少
$\pi\cdot1800^2/(0.866\cdot958^2)=12.8\approx13$ 个站位。

**这 13 个站位的位置由"全域覆盖"决定，与源在哪里无关**，
所以它们无法被"顺路经过源"替代。

\begin{center}
\begin{tabular}{lccc}
\toprule
格点间距 & 站位数 & 最差完成率 & 时间/s \\ \midrule
950（现行） & 24.4 & 0.9978 & 597.0 \\
1300 & 21.6 & 1.0000 & 646.4 \\
1500 & 22.7 & 1.0000 & 647.9 \\
1700 & 23.8 & 1.0000 & 694.4 \\
\bottomrule
\end{tabular}
\end{center}

\textbf{格点变疏并没有减少站位数}——因为认证搜索会自动把缺的补回来（自适应补站），
而且补出来的点路线更差，时间反而升高。

\subsection{减贴边环：行程几乎不变}

\begin{center}
\begin{tabular}{lcccc}
\toprule
贴边环点数 & 总站位数 & 最差完成率（本池） & 时间/s & 行程/m \\ \midrule
12 & 24.4 & 1.0000 & 582.2 & 23323 \\
6 & 18.9 & 1.0000 & 549.3 & 22477 \\
2 & 16.0 & 0.9928 & 566.4 & 23296 \\
0 & 14.9 & 0.9897 & 559.0 & 23498 \\
\bottomrule
\end{tabular}
\end{center}

\textbf{站位数从 24.4 减到 14.9（$-39\%$），行程只从 23323 减到 23498 m（$+0.7\%$）}。
原因是：贴边环被去掉后，认证搜索会为同样的边界候选点补出别的站位，
而补出来的点路线更差。\textbf{行程不是站位数造成的。}

\subsection{减归航中间点：会牺牲完成率}

每个源的清除需要若干次归航迭代（每次一次测量 + 一段移动）。
把定位门槛放宽（只在估计足够好时才去清）可以减少迭代，但会推迟清除、
增加"看见却清不掉"的风险。实测 8 组参数（定位门槛 250/300/400、
归航上限 260/420、末段半径 60/90、relocate 门槛 700/900、穷举触发 1/2）：

\begin{center}
\begin{tabular}{lcc}
\toprule
参数组 & 时间/s & 检测次数 \\ \midrule
基线 & 586.9 & 435 \\
最好（末段半径 60） & 581.1 & 438 \\
最差（穷举触发 1） & 597.8 & 434 \\
\bottomrule
\end{tabular}
\end{center}

\textbf{全部落在噪声内}，检测次数稳定在 434$\sim$440。说明这套流程的开销已经压到它自身的下界。

% ===========================================================================
\section{被否掉的路线（含结构性的重大改动）}

\subsection{让清除巡回兼任普查（自适应采样）}

\textbf{动机}：清除巡回本来就要走 8635 m 穿越全城，
若每 250 m 测量一次就是 35 个采样点——认证所需的读数很可能"免费"得到。

\textbf{实测}：

\begin{center}
\begin{tabular}{lcccc}
\toprule
配置 & 站位数 & 最差完成率 & 时间/s & 行程/m \\ \midrule
格点 $s{=}950$ $+$ 环12（现行） & 24.4 & 1.0000 & 582.2 & 23323 \\
自适应 batch4 & 24.9 & 0.9792 & 570.0 & 24558 \\
自适应 batch8 & 35.0 & 0.9603 & 520.2 & 21494 \\
自适应 batch12 & 38.6 & 0.9561 & 486.1 & 19704 \\
\bottomrule
\end{tabular}
\end{center}

自适应确实更快（486 s），但\textbf{完成率掉到 0.956}。原因是\textbf{自举失败}：
入场普查只听到少量源 $\to$ 清除巡回很短 $\to$ 沿途测量覆盖不足 $\to$ 发现更少 $\to$ 巡回更短。
巡回长度与发现率互相拖累，而且\textbf{贪心补站没有全局视野}，
一批一批地加导致路线反复穿越全域。

\subsection{粗格点保底 $+$ 沿途采样}

\textbf{动机}：全向源的发现只需覆盖半径 $\le R_{\min}=1000$ m（间距 $\le1732$ m，约 7 站），
认证交给沿途测量。

\textbf{实测}：完成率确实都是 1.0000，但\textbf{更慢}（646$\sim$694 s），
因为认证仍在补站、路线更乱。行程从 23323 m 升到 27381 m。

\subsection{环组布局（离线枚举优化）}

用旋转对称的环组（如"中心 $+$ 半径 900 的 6 点环 $+$ 半径 1650 的 12 点环"）
替代三角格点。离线精确评估 15 种组合，\textbf{最多只认证 88.6\% 的候选点}。

原因是环上的点在角度上太稀疏：位于两环之间的候选点，
其 $R$ 邻域内的读数集中在很窄的角度范围内，无法环绕它。
\textbf{三角格点才是认证的正确结构。}

顺带发现：\textbf{贪心集合覆盖在这个问题上无法启动}——
认证一个候选点需要至少 3 个环绕读数，单个站位的"边际增益"恒为 0。
这说明\textbf{认证不是"覆盖"问题，而是"环绕"问题}。

\subsection{全序列 Or-opt 与重规划策略}

把 Or-opt 从"只优化前 12 个任务"改成"O(1) 增量的全序列版本"——
\textbf{结果逐位相同}（605.3 s 对 605.3 s，行程 23441 m 对 23441 m）。
印证了第 1 节的结论：\textbf{巡回已经近最优}。

把"每个任务后重规划"改成"批量承诺既定巡回"——\textbf{更慢}
（605.3 $\to$ 643.0 s，行程 23441 $\to$ 25267 m），因为无法及时响应新发现的源。

% ===========================================================================
\section{本轮仍然取得的改进}

\begin{center}
\begin{tabular}{lcccc}
\toprule
 & 最差完成率 & 平均时间/s & 平均行程/m & 失败算例 \\ \midrule
优化前（P4） & 0.9978 & 880.7 & 41409 & 1/120 \\
\textbf{优化后} & \textbf{0.9978} & \textbf{597.0} & \textbf{33889} & 1/120 \\
\bottomrule
\end{tabular}
\end{center}

\textbf{完成率逐例相同，时间 $-32\%$。} 正式测试：问题四三局 14/14、12/12、14/14 全部清除，
平均 \textbf{579.2 s}（此前 1090.6 s，$-47\%$）。

改进来自三处：间距 1100$\to$950（认证可由格点自身完成）、
贴边环 12 点 @1750 m（把昂贵且难认证的边界候选点一次覆盖）、
认证边距 90$\to$700（只认证 $r\le1100$，边界发现交给贴边环）。
其中\textbf{贴边环是决定性的}：它替代掉的是自适应补站那种"为一两个候选点专门跑一趟"的行程。

% ===========================================================================
\section{若要真正逼近 300 s，需要改变问题的哪一项}

按第 1 节的分解，唯一的出路是\textbf{把"必须访问的点数"从 55 降到 20 以下}。
在现有设定下这需要以下任意一项成立：

\begin{enumerate}[leftmargin=1.8em]
  \item \textbf{认证判据被削弱}：例如允许"以一定置信度而非严格证明"停止搜索。
        代价是完成率不再有保证；
  \item \textbf{接收半径更大}：若 $R_{\min}$ 从 1000 m 提高到 1800 m，
        覆盖半径可放大到 $\sqrt3\times1800=3117$ m，认证只需 3$\sim$4 个站位，
        总点数可降到 25 左右，时间有望进入 350$\sim$400 s；
  \item \textbf{一次测量多个频道}：若测量动作能一次覆盖多个频道（而非每频道 5 s），
        435 次检测的 2610 s 可以大幅压缩。当前协议是每频道一次 \texttt{measure}，
        这项不在选手可控范围内。
\end{enumerate}

\begin{keybox}[title=一句话总结]
\textbf{不是算法不够好，而是"要走的路"本身就要 305 s/源。}
路由余量 8\%、检测次数已在自身下界、站位数与行程无关——
这三条实测堵住了继续压缩的通道。因此本文把 200$\sim$300 s 记为
\textbf{在本问题设定下不可达}，并保留 597 s $/$ 0.9978 作为交付配置。
\end{keybox}
\end{document}
